\documentclass{subfiles}
\begin{document}
    Taking the name from the three researchers who developed it, \citeauthor{RivestSA1977},
        RSA is one of the simplest, yet solid\footnotemark algorithm in cryptography.

    The idea behind RSA is the following: we choose two prime numbers \(p\) and \(q\)
        of similar size, i.e., their binary representation uses approximately the same numebr of bits,
        we compute their product \(n\) and use this value to cypher our message. 
        In \autoref{proc:RSA-schema}, we give a more detailed description of the RSA algorithm.

        \subfile{../TikZ/Procedure - RSA schema}

    \footnotetext{
        Here by ``solid'' we refer to the fact that cracking RSA is somewhat unfeasable.
    }

    What we are left with is showing that Line 3 of \autoref{proc:RSA-schema},
        does in fact return \(m\). First observe that 
        \[
            c^{d} = (m^{e}) = m^{ed}.
        \]
        Using Bézout identity, we get 
        \[
            m^{ed} = m^{1 - \mu\lambda} = m * m^{-\mu\lambda} = m * (m^{\lambda})^{-\mu}.
        \]
    
    Recall that by definition \(\lambda = \lcm(p - 1, q - 1)\),
        hence there exists \(r, s\) such that \(\lambda = (p -1)r = (q - 1)s\).
        By using Fermat's theorem, which roughuly states that \(x^{p} \equiv x \mod\),
        we get
        \[
            m^{\lambda} = m^{(p -1)r} = (m^{(p - 1)})^{r} = 1^{r} = 1 \mod.
        \]
        By the same logic, we get that \(m^{\lambda} = 1 \mod[q]\).
        Therefore,
        \[
            c^{d} = m^{ed} = (m^{\lambda})^{-\mu} = m * 1 \mod = m * 1 \mod[q] = m.
        \]
\end{document}
